\documentclass{article}
\usepackage{amsmath}
\usepackage{listings}
\usepackage{microtype}
\usepackage{tikz}
\usepackage{pgfplots}
\usepackage{booktabs}
\usetikzlibrary{arrows.meta,positioning}
\usepackage[ngerman]{babel}
\begin{document}

\section{Cluster-Infrastruktur und Parallelisierung der Ray-Tracing-Engine}
\label{sec:cluster}

\subsection{Aufgabenstellung}
Im Rahmen des Projekts bestand die Aufgabe darin, die Rust-basierte Physics
Engine auf einem Cluster aus acht Worker-Nodes parallel auszuführen und die Teilergebnisse zu 
einer konsistenten Impulsantwort (IR) zusammenzuführen. Dies umfasste (1) die Einrichtung einer 
passwortlosen SSH-Infrastruktur zwischen Manager- und Worker-Nodes, (2) die Kompilierung und 
Verteilung der Binary sowie (3) die
Entwicklung von Automatisierungs-Skripten für Ausführung und Merge.

\subsection{Architektur}
Abbildung~\ref{fig:cluster-arch} zeigt die eingesetzte Manager-Worker-Topologie.
Der Manager-Node orchestriert die Ausführung über SSH und sammelt die
Ergebnisdateien nach Abschluss der Simulation ein.

\begin{figure}[htbp]
 \centering
 \begin{tikzpicture}[
 node distance=0.5cm and 0.5cm,
 every node/.style={font=\small},
 manager/.style={rectangle, rounded corners, draw=black, thick,
 fill=blue!10, minimum width=1.5cm, minimum height=1cm, align=center},
 worker/.style={rectangle, rounded corners, draw=black,
 fill=gray!10, minimum width=0.5cm, minimum height=0.85cm, align=center},
 every edge/.style={draw, -{Latex[length=2mm]}, thin}
 ]
 \node[manager] (mgr) {Manager-Node\\ \texttt{run\_cluster.sh}};
 \node[worker, below left=1.0cm and 3.2cm of mgr] (n1) {Node 1\\ \texttt{rust\_service}};
 \node[worker, right=0.3cm of n1] (n2) {Node 2};
 \node[worker, right=0.3cm of n2] (n3) {Node 3};
 \node[worker, right=0.3cm of n3] (n4) {Node 4};
 \node[worker, right=0.3cm of n4] (n5) {Node 5};
 \node[worker, right=0.3cm of n5] (n6) {Node 6};
 \node[worker, right=0.3cm of n6] (n7) {Node 7};
 \node[worker, right=0.3cm of n7] (n8) {Node 8};
 \foreach \n in {n1,n2,n3,n4,n5,n6,n7,n8}{
 \draw[-{Latex[length=1.0mm]}] (mgr.south) -- (\n.north)
 node[midway, sloped, above, font=\tiny] {};
 }
 \node[below=0.15cm of n1, font=\tiny, align=center] {1M Rays};
 \node[below=0.15cm of n8, font=\tiny, align=center] {1M Rays};
 \end{tikzpicture}
 \caption{Manager-Worker-Topologie: Der Manager verteilt via SSH die
 kompilierte Binary sowie \texttt{input\_config.json} und startet die
 Ausführung auf allen acht Worker-Nodes parallel (Hintergrundprozesse,
 synchronisiert über \texttt{wait}).}
 \label{fig:cluster-arch}
\end{figure}

\subsection{Implementierung}

\paragraph{SSH-Infrastruktur.}
Für die automatisierte, nicht-interaktive Ausführung wurde ein
Ed25519-Schlüsselpaar auf dem Manager-Node erzeugt und mittels
\texttt{ssh-copy-id} auf alle acht Worker-Nodes verteilt. Die
Knotenadressierung erfolgte über statische IP-Adressen anstelle von
Hostnamen, da im Cluster-Netzwerk keine Namensauflösung konfiguriert war.

\paragraph{Build \& Deployment.}
Die Rust-Binary wurde zentral auf dem Manager-Node aus dem Unterverzeichnis
\texttt{rust\_service} kompiliert (\texttt{cargo build --release},
Rust-Toolchain via \texttt{rustup}, Edition 2024) und anschließend
zusammen mit \texttt{input\_config.json} per \texttt{scp} auf alle Nodes
verteilt.

\paragraph{Orchestrierung.}
Das Skript \texttt{run\_cluster.sh} startet die Binary auf allen acht
Nodes parallel als Hintergrundprozesse und synchronisiert deren Abschluss
über \texttt{wait}. Nach Terminierung aller Prozesse werden die acht
Ergebnisdateien (\texttt{ir\_node1.json}~--~\texttt{ir\_node8.json}) auf
den Manager zurückkopiert.

\paragraph{Merge.}
Das Python-Skript \texttt{merge\_results.py} konsolidiert die Felder
\texttt{delays\_seconds} und \texttt{pressures} aller acht Teilergebnisse
zu einer einheitlichen Impulsantwort \texttt{ir\_merged.json}, welche
anschließend an die DSP-Pipeline übergeben wird.

\subsection{Theoretische Grundlagen: Speedup und Amdahls Gesetz}
\label{sec:amdahl}
Um die im Folgenden gemessenen Laufzeiten einordnen zu
können, wird zunächst der theoretische Rahmen zur Bewertung paralleler
Systeme skizziert.

\paragraph{Speedup.}
Der Speedup $S(n)$ bei Verwendung von $n$ parallelen Verarbeitungseinheiten
(hier: Nodes) ist definiert als das Verhältnis der seriellen Laufzeit
$T(1)$ zur parallelen Laufzeit $T(n)$:
\begin{equation}
 S(n) = \frac{T(1)}{T(n)}
 \label{eq:speedup}
\end{equation}
Ein Speedup von $S(n) = n$ wird als \emph{linear} bezeichnet und stellt
den theoretischen Idealfall dar, bei dem eine Verdopplung der Ressourcen
die Laufzeit exakt halbiert.

\paragraph{Amdahls Gesetz.}
In der Praxis lässt sich jedoch nicht jeder Programmteil parallelisieren.
Amdahls Gesetz\footnote{Amdahl, G. M. (1967). \emph{Validity of the
single processor approach to achieving large scale computing
capabilities}. AFIPS Conference Proceedings, 30, 483--485. --
Sobald eine gemeinsame \texttt{.bib}-Datei im Projekt vorliegt, sollte
dieser Verweis durch \texttt{\textbackslash cite\{amdahl1967\}}
ersetzt werden.} beschreibt den maximal erreichbaren
Speedup in Abhängigkeit vom parallelisierbaren Anteil $p$ eines Programms
(bei einem seriellen, nicht parallelisierbaren Restanteil von $1-p$):
\begin{equation}
 S(n) = \frac{1}{(1-p) + \dfrac{p}{n}}
 \label{eq:amdahl}
\end{equation}
Für $n \to \infty$ konvergiert der Speedup gegen den Grenzwert
$\frac{1}{1-p}$. Das bedeutet: Selbst bei einer beliebig großen Anzahl an
Nodes limitiert bereits ein kleiner serieller Anteil den maximal
erreichbaren Gesamtspeedup. Ein Programm mit z.\,B. 10\,\% seriellem
Anteil ($p = 0{,}9$) kann demnach nie schneller als das 10-fache der
Ausgangslaufzeit werden, unabhängig von der Anzahl eingesetzter Nodes.

\paragraph{Anwendung auf das Ray-Tracing-Problem.}
Das hier vorliegende Simulationsproblem gehört zur Klasse der
\emph{embarrassingly parallel} Probleme: Jeder der $N_{\text{total}}$
Schallstrahlen kann vollständig unabhängig von allen anderen berechnet
werden, da während der Ray-Generierung kein Datenaustausch zwischen den
Nodes erforderlich ist. Die Aufteilung erfolgt daher denkbar einfach durch
gleichmäßige Verteilung der Gesamt-Ray-Anzahl auf die $n$ Worker-Nodes
($N_{\text{total}} / n$ Rays pro Node), analog zur Verteilung durch
\texttt{rayon} innerhalb eines einzelnen Nodes.

Der serielle Anteil $1-p$ beschränkt sich in diesem Aufbau im Wesentlichen
auf die Orchestrierungsphasen vor und nach der eigentlichen Berechnung:
den Verbindungsaufbau via SSH, die Verteilung von Binary und
Konfigurationsdatei per \texttt{scp}, sowie das abschließende Einsammeln
und Zusammenführen (\texttt{merge\_results.py}) der Ergebnisdateien. Diese
Schritte sind nicht parallelisierbar, machen jedoch bei einer
Rechenlast von 1\,Million Rays pro Node und einer Laufzeit von wenigen
Millisekunden pro Node (siehe Tabelle~\ref{tab:results}) nur einen
verschwindend geringen Anteil an der Gesamtlaufzeit aus. Der
parallelisierbare Anteil $p$ liegt damit sehr nahe an $1$, wodurch nach
Gleichung~\eqref{eq:amdahl} ein nahezu linearer Speedup zu erwarten ist
-- eine Erwartung, die sich mit dem in Abschnitt~\ref{sec:finaler-testlauf}
dokumentierten Testlauf deckt, in dem sich die Rechenzeit pro Node bei
$n=8$ gegenüber $n=1$ praktisch nicht verändert (vgl.\ Abbildung~\ref{fig:durchsatz}).

\paragraph{Grenzen der Skalierung.}
Zu beachten ist, dass der serielle Anteil mit wachsender Node-Anzahl nicht
konstant bleibt: Verbindungsaufbau und Dateiübertragung 
pro Node
verursachen einen Kommunikationsoverhead, der mit $n$ wächst
(bei $n$ Nodes sind $n$ SSH-Verbindungen und $n$ \texttt{scp}-Transfers
erforderlich). Bei den hier verwendeten $n=8$ Nodes ist dieser Effekt noch
nicht sichtbar; bei einer deutlich größeren Anzahl an Nodes würde jedoch
zu erwarten sein, dass der gemessene Speedup zunehmend vom linearen
Idealfall abweicht, da der serielle Anteil $1-p$ effektiv ansteigt. Diese
Überlegung liefert die theoretische Grundlage für den unten skizzierten
Ausblick auf eine dynamische, lastabhängige Verteilung.

\subsection{Ergebnisse}
\label{sec:ergebnisse}
\label{sec:finaler-testlauf}
Der initiale Testlauf mit einer stark reduzierten Ray-Anzahl
(\texttt{rays\_to\_cast = 1}) diente lediglich zur Verifikation der
grundsätzlichen Funktionsfähigkeit der Pipeline. Für einen
aussagekräftigen Vergleich der Ausführung auf einem einzelnen Node
gegenüber der verteilten Ausführung auf acht Nodes wurde daher ein
weiterer Testlauf mit einer produktionsnahen Konfiguration
(\texttt{rays\_to\_cast = 100000}, \texttt{mic\_radius = 0.5}) auf dem
Branch \texttt{3d-audio-cluster} durchgeführt.

Bei der Auswertung ist zwischen zwei unterschiedlichen Zeitmessungen zu
unterscheiden: der \emph{Wandzeit} (\texttt{real}-Zeit des
Unix-Kommandos \texttt{time}), die neben der reinen
Simulationsberechnung auch den SSH-Verbindungsaufbau sowie den
abschließenden Rücktransfer der acht Ergebnisdateien per \texttt{scp}
einschließt, und der von der Rust-Engine selbst intern gemessenen
\emph{Simulationslaufzeit} (\texttt{Simulation run time}), welche
ausschließlich die reine Berechnungszeit ohne Kommunikations-Overhead
widerspiegelt.

Tabelle~\ref{tab:finaler-testlauf} stellt die gemessenen Werte für die
Ausführung auf einem einzelnen Node ($n=1$) der Ausführung auf acht
Nodes ($n=8$) gegenüber.

\begin{table}[htbp]
 \centering
 \caption{Vergleich Einzel-Node- vs. Acht-Node-Ausführung bei
 \texttt{rays\_to\_cast = 100000} pro Node}
 
 \label{tab:finaler-testlauf}
 \begin{tabular}{lrr}
 \toprule
 Metrik & $n=1$ & $n=8$ \\
 \midrule
 Rays pro Node & 100.000 & 100.000 \\
 Rays gesamt & 100.000 & 800.000 \\
 Simulationslaufzeit pro Node (intern) & 0{,}044\,s & 0{,}033--0{,}044\,s \\
 Wandzeit (\texttt{real}, gesamter Lauf) & 1{,}275\,s & 3{,}813\,s \\
 Ray-Hits gesamt (Merge) & 141.115 & 1.133.347 \\
 \bottomrule
 \end{tabular}
\end{table}

Die reine, intern gemessene Simulationslaufzeit liegt sowohl beim
Einzel-Node- als auch bei jedem der acht parallel laufenden Nodes im
selben engen Bereich von rund $0{,}033$ bis $0{,}044\,$s -- unabhängig
davon, ob ein Node isoliert oder gemeinsam mit sieben weiteren Nodes
gleichzeitig rechnet. Die insgesamt verarbeitete Ray-Anzahl steigt bei
gleichbleibender Rechenzeit pro Node hingegen linear mit der Anzahl der
Nodes: Bei $n=8$ werden in etwa derselben Zeitspanne, die ein einzelner
Node für $100.000$ Rays benötigt, insgesamt $800.000$ Rays verarbeitet.

Demgegenüber steigt die gemessene Wandzeit des Gesamtlaufs
(\texttt{real}) von $1{,}275\,$s bei $n=1$ auf $3{,}813\,$s bei $n=8$ an.
Dieser Anstieg ist nicht auf die Simulationsberechnung selbst
zurückzuführen, sondern resultiert vollständig aus dem seriellen, nicht
parallelisierten Anteil des Ablaufs: dem Aufbau von acht
SSH-Verbindungen sowie insbesondere dem sequenziellen Rücktransfer der
acht ca.\ $4\,$MB großen \texttt{ir\_output.json}-Dateien per
\texttt{scp} nach Abschluss der Berechnung. Genau dieser Anteil
entspricht dem seriellen Restanteil $1-p$ aus
Gleichung~\eqref{eq:amdahl} in Abschnitt~\ref{sec:amdahl} und bestätigt
empirisch die dort formulierte Erwartung, dass der
Kommunikations-Overhead mit wachsender Node-Anzahl $n$ zunimmt.

Der vorliegende Testlauf zeigt damit kein klassisches
Latenz-Speedup im Sinne von $S(n) = T(1)/T(n)$ mit gleichbleibender
Gesamt-Ray-Anzahl -- ein solcher Vergleich wäre hier irreführend, da bei
$n=1$ und $n=8$ unterschiedliche Gesamt-Ray-Mengen berechnet wurden.
Stattdessen zeigt sich eine \emph{Durchsatzskalierung}
(Abbildung~\ref{fig:durchsatz}): Bei nahezu konstanter Rechenzeit pro
Node lässt sich die Gesamtzahl der verarbeiteten Rays proportional zur
Anzahl der eingesetzten Nodes steigern -- die für ein \emph{embarrassingly
parallel} Problem (vgl.\ Abschnitt~\ref{sec:amdahl}) zu erwartende
Charakteristik.

\begin{figure}[htbp]
 \centering
 \begin{tikzpicture}
 \begin{axis}[
 width=0.72\textwidth,
 height=6.5cm,
 ybar,
 bar width=25pt,
 symbolic x coords={n=1, n=8},
 xtick=data,
 xlabel={Anzahl Nodes},
 ylabel={Verarbeitete Rays gesamt},
 ymin=0,
 grid=major,
 grid style={dashed, gray!30},
 nodes near coords,
 nodes near coords align={vertical},
 every node near coord/.append style={font=\small},
 legend pos=north west,
 legend style={font=\small}
 ]
 \addplot[fill=blue!40, draw=blue] coordinates {
 (n=1,100000) (n=8,800000)
 };
 \addlegendentry{Rays gesamt (100.000 pro Node)}
 \end{axis}
 \end{tikzpicture}
 \caption{Durchsatzskalierung der verteilten Ray-Tracing-Berechnung: Bei
 nahezu konstanter Simulationslaufzeit pro Node
 ($0{,}033$--$0{,}044\,$s, vgl.\ Tabelle~\ref{tab:finaler-testlauf})
 steigt die insgesamt verarbeitete Ray-Anzahl proportional zur Anzahl
 der eingesetzten Nodes von $100.000$ ($n=1$) auf $800.000$ ($n=8$).
 Im Gegensatz zu einem klassischen Latenz-Speedup-Diagramm zeigt diese
 Darstellung die für ein \emph{embarrassingly parallel} Problem
 charakteristische lineare Durchsatzsteigerung bei konstanter
 Rechenzeit.}
 \label{fig:durchsatz}
\end{figure}

\subsection{Herausforderungen und Lösungen}
Die technischen Schwierigkeiten lagen primär in der Infrastruktur und
nicht in der Simulationslogik selbst; sie betrafen im Wesentlichen die
Etablierung eines zuverlässigen, nicht-interaktiven Zugriffs auf die
Worker-Nodes.

Den größten Zeitaufwand verursachte dabei bereits die Grundvoraussetzung
für jede weitere Konfiguration: der Zugriff auf mehrere Worker-Nodes selbst.
Da für diese keine funktionierenden Zugangsdaten vorlagen, mussten die
Passwörter zunächst über ein Live-USB-Boot-Medium und anschließenden
chroot-Zugriff auf das jeweilige Root-Filesystem zurückgesetzt werden --
ein Schritt, der der eigentlichen SSH-Schlüsselverteilung vorausgehen
musste und der Infrastrukturarbeit im Projekt somit ein erhebliches
Gewicht verlieh, das sich im späteren Code nicht mehr abbildet.

Nach Wiederherstellung des Zugriffs zeigte sich, dass das initial
erzeugte SSH-Schlüsselpaar versehentlich auf dem falschen Node erstellt
worden war; mittels \texttt{ssh -v}-Diagnose ließ sich diese
Fehlkonfiguration identifizieren und korrigieren. Ein verwandtes Problem
ergab sich bei der Adressierung der Nodes: Da im Cluster-Netzwerk keine
Namensauflösung konfiguriert war, scheiterten hostnamenbasierte
SSH-Verbindungen durchweg, sodass stattdessen auf eine direkte
Adressierung über die mittels \texttt{hostname -I} ermittelten statischen
IP-Adressen zurückgegriffen wurde.

Erst nachdem SSH-Zugriff und Adressierung zuverlässig funktionierten,
trat mit der Toolchain-Kompatibilität eine Schwierigkeit anderer Art
zutage: Die auf den Nodes vorinstallierte Cargo-Version (1.75)
unterstützte die von der Physics Engine verwendete \texttt{edition2024}
nicht, was eine Neuinstallation der Rust-Toolchain via \texttt{rustup}
auf allen Nodes erforderlich machte.

Insgesamt zeigt sich an diesen vier Problemen ein gemeinsames Muster:
Nahezu der gesamte Aufwand entfiel auf die Herstellung eines
funktionierenden Ausgangszustands der Infrastruktur, während die
eigentliche Parallelisierungslogik in \texttt{run\_cluster.sh} vergleichsweise
unproblematisch umzusetzen war.

\subsection{Ausblick}
\label{sec:ausblick}
Als Weiterentwicklung bietet sich eine \emph{dynamische Lastverteilung}
an, bei der die Ray-Anzahl pro Node anhand der individuellen
Rechenleistung skaliert wird, sowie ein automatisiertes Monitoring der
Node-Auslastung zur weiteren Effizienzsteigerung.

\end{document}
